% Generated draft: requires mathematical and source review.
\documentclass[11pt]{article}
\usepackage[margin=25mm]{geometry}
\usepackage{fontspec}
\setmainfont{Latin Modern Roman}
\usepackage{amsmath,amssymb,mathtools,unicode-math}
\setmathfont{Latin Modern Math}
\usepackage{enumitem,needspace,xurl,hyperref}
\setlength{\parindent}{0pt}
\setlength{\parskip}{5pt}
\setlength{\emergencystretch}{3em}
\newcommand{\N}{\mathbb{N}}
\newcommand{\Z}{\mathbb{Z}}
\newcommand{\Q}{\mathbb{Q}}
\newcommand{\R}{\mathbb{R}}
\newcommand{\C}{\mathbb{C}}
\newcommand{\F}{\mathbb{F}}
\newcommand{\A}{\mathbb{A}}
\newcommand{\PP}{\mathbb{P}}
\newcommand{\E}{\mathbb{E}}
\newcommand{\eps}{\varepsilon}
\newcommand{\dd}{\,\mathrm{d}}
\begin{document}
\section*{UnsolvedMath ID: 700878}
\textbf{Source problem code:} catalogue entry
\subsection{Definitions and mathematical statement}
A strongly regular graph with parameters \((85,14,3,2)\) means a finite simple undirected graph \(G=(V,E)\) such that \(|V|=85\), every vertex has exactly \(14\) neighbors, every two distinct adjacent vertices have exactly \(3\) common neighbors, and every two distinct nonadjacent vertices have exactly \(2\) common neighbors. For \(v\in V\), write \(N_G(v)=\{u\in V:\{u,v\}\in E\}\); the neighborhood graph of \(v\) is the induced graph \(G[N_G(v)]\).

Does there exist such a graph \(G\) and a vertex \(v\in V\) for which \(G[N_G(v)]\) is disconnected? Explicitly, the additional requirement is that there exist nonempty disjoint sets \(A,B\subseteq N_G(v)\) with \(N_G(v)=A\cup B\) and \(\{a,b\}\notin E\) for every \(a\in A\) and \(b\in B\).
\subsection{Short English statement}
Does there exist a strongly regular graph with parameters \((85,14,3,2)\) in which the subgraph induced by the neighbors of at least one vertex is disconnected?
\subsection{Sources}
\begin{itemize}
\item Local UnsolvedMath registry: supplied record with permanent ID \(700878\), exact source problem code \texttt{catalogue entry}, titled \emph{Kourovka Notebook Problem 12.59}.
\item Supplied source attribution: \emph{The Kourovka Notebook (Unsolved Problems in Group Theory)}, original problem 12.59, 12th issue (1992); proposed by A. A. Makhn\"ev.
\item The supplied background also identifies the 21st issue edition (2026), arXiv:1401.0300v48. This bibliographic attribution is reproduced from the record and has not been independently verified.
\end{itemize}
\end{document}
